% GENERATED, edit the French unit or the translation table
\mxunit{stmg-tests}{Supervised tests}
Two versions are provided for each chapter. The reference structure is: fluency questions worth 6 points, an exercise worth 6 points and a problem worth 8 points.
\section*{Chapter 1 -- Numerical sequences}
\mxtest{1}{A}{Numerical sequences}{model, calculate, reason, communicate}
\begin{exercise}[1 -- Fluency]\mxpoints{6}
\begin{enumerate}
\item Give the multiplier associated with an increase of $5\,\%$.
\item Is the sequence $12,17,22,27,\ldots$ arithmetic? Give its common difference.
\item How many terms does $u_4+u_5+\cdots+u_{18}$ contain?
\item If $w_0=80$ and $w_{n+1}=0.9w_n$, calculate $w_2$.
\item Write $1+2+\cdots+20$ using the symbol $\sum$.
\end{enumerate}
\end{exercise}
\begin{exercise}[2 -- Subscriptions]\mxpoints{6}
A platform has $120\,000$ subscribers at the start of the year and gains $11\,000$ each month. Let $u_n$ be the number of subscribers, in thousands, after $n$ months.
\begin{enumerate}
\item Show that $(u_n)$ is arithmetic and state its common difference.
\item Express $u_n$ in terms of $n$, then calculate $u_{12}$.
\item Calculate $u_1+\cdots+u_{12}$. Interpret this sum with caution.
\end{enumerate}
\end{exercise}
\begin{exercise}[3 -- Investment and target]\mxpoints{8}
An initial capital of $650\,000$ Ar increases by $5\,\%$ per year. Let $C_n$ be the capital after $n$ years.
\begin{enumerate}
\item Write $C_{n+1}$ in terms of $C_n$, then $C_n$ in terms of $n$.
\item Calculate $C_5$, rounded to the nearest ariary.
\item Use a calculator to find the first integer $n$ for which $C_n$ exceeds $1\,000\,000$ Ar.
\item Explain why the model is geometric and not arithmetic.
\end{enumerate}
\end{exercise}
\begin{solution}\textbf{Marking scheme.} Ex. 1: 6 points spread over the fluency questions. Ex. 2: 6 points. Ex. 3: 8 points.\\
\textbf{Solution outline.} The multiplier is $1.05$. The first sequence in the fluency questions is arithmetic with common difference $5$; the sum contains $15$ terms; $w_2=64$ and $\sum_{k=1}^{20}k$. For exercise 2, $u_n=120+11n$, so $u_{12}=252$; the requested sum uses $12\dfrac{u_1+u_{12}}2$ and accumulates monthly figures, not a number of distinct people. For exercise 3, $C_n=650\,000\times\left(1+\dfrac{5}{100}\right)^n$; check the threshold with two consecutive values.\end{solution}

\mxtest{1}{B}{Numerical sequences}{model, calculate, reason, communicate}
\begin{exercise}[1 -- Fluency]\mxpoints{6}
\begin{enumerate}
\item Give the multiplier associated with an increase of $4\,\%$.
\item Is the sequence $12,17,22,27,\ldots$ arithmetic? Give its common difference.
\item How many terms does $u_4+u_5+\cdots+u_{18}$ contain?
\item If $w_0=80$ and $w_{n+1}=0.9w_n$, calculate $w_2$.
\item Write $1+2+\cdots+20$ using the symbol $\sum$.
\end{enumerate}
\end{exercise}
\begin{exercise}[2 -- Subscriptions]\mxpoints{6}
A platform has $145\,000$ subscribers at the start of the year and gains $9\,000$ each month. Let $u_n$ be the number of subscribers, in thousands, after $n$ months.
\begin{enumerate}
\item Show that $(u_n)$ is arithmetic and state its common difference.
\item Express $u_n$ in terms of $n$, then calculate $u_{12}$.
\item Calculate $u_1+\cdots+u_{12}$. Interpret this sum with caution.
\end{enumerate}
\end{exercise}
\begin{exercise}[3 -- Investment and target]\mxpoints{8}
An initial capital of $720\,000$ Ar increases by $4\,\%$ per year. Let $C_n$ be the capital after $n$ years.
\begin{enumerate}
\item Write $C_{n+1}$ in terms of $C_n$, then $C_n$ in terms of $n$.
\item Calculate $C_5$, rounded to the nearest ariary.
\item Use a calculator to find the first integer $n$ for which $C_n$ exceeds $1\,000\,000$ Ar.
\item Explain why the model is geometric and not arithmetic.
\end{enumerate}
\end{exercise}
\begin{solution}\textbf{Marking scheme.} Ex. 1: 6 points spread over the fluency questions. Ex. 2: 6 points. Ex. 3: 8 points.\\
\textbf{Solution outline.} The multiplier is $1.04$. The first sequence in the fluency questions is arithmetic with common difference $5$; the sum contains $15$ terms; $w_2=64$ and $\sum_{k=1}^{20}k$. For exercise 2, $u_n=145+9n$, so $u_{12}=253$; the requested sum uses $12\dfrac{u_1+u_{12}}2$ and accumulates monthly figures, not a number of distinct people. For exercise 3, $C_n=720\,000\times\left(1+\dfrac{4}{100}\right)^n$; check the threshold with two consecutive values.\end{solution}

\section*{Chapter 2 -- Exponential functions}
\mxtest{2}{A}{Exponential functions}{calculate, represent, model, interpret}
\begin{exercise}[1 -- Powers and variations]\mxpoints{6}
\begin{enumerate}
\item Simplify $2^4\times2^{-7}$ and $\dfrac{5^{x+3}}{5^x}$.
\item Write $3^{-2}$ as a fraction.
\item Give the direction of variation of $f(x)=1.04^x$ and of $g(x)=-2\times0.75^x$.
\item Explain why $a^x+a^y=a^{x+y}$ is false in general.
\end{enumerate}
\end{exercise}
\begin{exercise}[2 -- Website traffic]\mxpoints{6}
At launch, a website receives $2400$ visits per day. Traffic increases by $4\,\%$ per week. It is modelled by $V(t)=2400\times1.04^t$.
\begin{enumerate}
\item Interpret the numbers $2400$ and $1.04$.
\item Calculate $V(8)$, rounded to the nearest unit.
\item Find the first whole number of weeks for which $V(t)>3100$.
\end{enumerate}
\end{exercise}
\begin{exercise}[3 -- Average rate]\mxpoints{8}
A company's turnover goes from $82\,000$ to $104000$ euros in four years.
\begin{enumerate}
\item Calculate the overall multiplier.
\item Write the average annual multiplier as a power.
\item Calculate the average annual rate of change (percentage change), rounded to $0.01\,\%$.
\item Explain why the overall rate is not simply divided by four.
\end{enumerate}
\end{exercise}
\begin{solution}\textbf{Marking scheme:} $6+6+8$. \textbf{Solution:} $2^4\times2^{-7}=2^{-3}=\dfrac18$, $\dfrac{5^{x+3}}{5^x}=125$, $3^{-2}=\dfrac19$. $f$ is increasing and $g$ is increasing. In the problem, calculate $V(8)$ then test two consecutive integers around the threshold. For the average rate, $q=\left(\dfrac{V_f}{V_i}\right)^{\frac14}$ and $t=q-1$; multipliers multiply.\end{solution}

\mxtest{2}{B}{Exponential functions}{calculate, represent, model, interpret}
\begin{exercise}[1 -- Powers and variations]\mxpoints{6}
\begin{enumerate}
\item Simplify $2^4\times2^{-7}$ and $\dfrac{5^{x+3}}{5^x}$.
\item Write $3^{-2}$ as a fraction.
\item Give the direction of variation of $f(x)=1.03^x$ and of $g(x)=-2\times0.75^x$.
\item Explain why $a^x+a^y=a^{x+y}$ is false in general.
\end{enumerate}
\end{exercise}
\begin{exercise}[2 -- Website traffic]\mxpoints{6}
At launch, a website receives $1850$ visits per day. Traffic increases by $3\,\%$ per week. It is modelled by $V(t)=1850\times1.03^t$.
\begin{enumerate}
\item Interpret the numbers $1850$ and $1.03$.
\item Calculate $V(8)$, rounded to the nearest unit.
\item Find the first whole number of weeks for which $V(t)>2250$.
\end{enumerate}
\end{exercise}
\begin{exercise}[3 -- Average rate]\mxpoints{8}
A company's turnover goes from $82\,000$ to $98000$ euros in four years.
\begin{enumerate}
\item Calculate the overall multiplier.
\item Write the average annual multiplier as a power.
\item Calculate the average annual rate of change (percentage change), rounded to $0.01\,\%$.
\item Explain why the overall rate is not simply divided by four.
\end{enumerate}
\end{exercise}
\begin{solution}\textbf{Marking scheme:} $6+6+8$. \textbf{Solution:} $2^4\times2^{-7}=2^{-3}=\dfrac18$, $\dfrac{5^{x+3}}{5^x}=125$, $3^{-2}=\dfrac19$. $f$ is increasing and $g$ is increasing. In the problem, calculate $V(8)$ then test two consecutive integers around the threshold. For the average rate, $q=\left(\dfrac{V_f}{V_i}\right)^{\frac14}$ and $t=q-1$; multipliers multiply.\end{solution}

\section*{Chapter 3 -- The common logarithm function}
\mxtest{3}{A}{Common logarithm}{calculate, solve, reason, communicate}
\begin{exercise}[1 -- Knowledge and calculations]\mxpoints{6}
\begin{enumerate}
\item Calculate $\log(1000)$, $\log(0.01)$ and $\log(10^{-4})$.
\item For $x,y>0$, expand $\log\!\left(\dfrac{100x^3}{y} \right)$.
\item Solve $10^t=28000$ and give $t$ to the nearest hundredth.
\end{enumerate}
\end{exercise}
\begin{exercise}[2 -- Growth target]\mxpoints{6}
A value is multiplied by $1.05$ each year. We look for the time needed for it to double.
\begin{enumerate}
\item Translate the problem into an equation.
\item Solve this equation using the common logarithm (base 10).
\item Give the first whole number of years that reaches or exceeds double.
\item Check with the two relevant consecutive powers.
\end{enumerate}
\end{exercise}
\begin{exercise}[3 -- Sound level]\mxpoints{8}
The sound level is $L=10\log\!\left(\dfrac I{I_0}\right)$, in decibels.
\begin{enumerate}
\item Calculate $L$ when $I=10^6I_0$.
\item For $L=84$, express then calculate the ratio $\dfrac I{I_0}$.
\item Compare the intensities corresponding to $84$ dB and $74$ dB.
\end{enumerate}
\end{exercise}
\begin{solution}\textbf{Marking scheme:} $6+6+8$. \textbf{Solution:} The usual values are $3,-2,-4$. The expansion gives $2+3\log x-\log y$. $t=\log(28000)$. For the doubling, $x=\dfrac{\log2}{\log(1.05)}$ then take the next integer up and check. $\dfrac I{I_0}=10^{\frac L{10}}$; a difference of $10$ dB corresponds to an intensity multiplied by $10$.\end{solution}

\mxtest{3}{B}{Common logarithm}{calculate, solve, reason, communicate}
\begin{exercise}[1 -- Knowledge and calculations]\mxpoints{6}
\begin{enumerate}
\item Calculate $\log(1000)$, $\log(0.01)$ and $\log(10^{-4})$.
\item For $x,y>0$, expand $\log\!\left(\dfrac{100x^3}{y} \right)$.
\item Solve $10^t=32000$ and give $t$ to the nearest hundredth.
\end{enumerate}
\end{exercise}
\begin{exercise}[2 -- Growth target]\mxpoints{6}
A value is multiplied by $1.07$ each year. We look for the time needed for it to double.
\begin{enumerate}
\item Translate the problem into an equation.
\item Solve this equation using the common logarithm (base 10).
\item Give the first whole number of years that reaches or exceeds double.
\item Check with the two relevant consecutive powers.
\end{enumerate}
\end{exercise}
\begin{exercise}[3 -- Sound level]\mxpoints{8}
The sound level is $L=10\log\!\left(\dfrac I{I_0}\right)$, in decibels.
\begin{enumerate}
\item Calculate $L$ when $I=10^6I_0$.
\item For $L=78$, express then calculate the ratio $\dfrac I{I_0}$.
\item Compare the intensities corresponding to $78$ dB and $68$ dB.
\end{enumerate}
\end{exercise}
\begin{solution}\textbf{Marking scheme:} $6+6+8$. \textbf{Solution:} The usual values are $3,-2,-4$. The expansion gives $2+3\log x-\log y$. $t=\log(32000)$. For the doubling, $x=\dfrac{\log2}{\log(1.07)}$ then take the next integer up and check. $\dfrac I{I_0}=10^{\frac L{10}}$; a difference of $10$ dB corresponds to an intensity multiplied by $10$.\end{solution}

\section*{Chapter 4 -- The reciprocal function}
\mxtest{4}{A}{Reciprocal function}{calculate, reason, represent, communicate}
\begin{exercise}[1 -- Reciprocal function]\mxpoints{6}
Consider $f(x)=2+\dfrac{16}x$.
\begin{enumerate}
\item Give its domain.
\item Calculate its derivative and study its sign.
\item Give its asymptotes and its direction of variation on each interval of the domain.
\end{enumerate}
\end{exercise}
\begin{exercise}[2 -- Sign table]\mxpoints{6}
Solve $\dfrac{(x-4)(x+2)}x\geq0$. The table must show the zeros and the excluded value.
\end{exercise}
\begin{exercise}[3 -- Average cost]\mxpoints{8}
For $x>0$, an average cost is modelled by $C(x)=x+\dfrac{16}x+12$.
\begin{enumerate}
\item Calculate $C'(x)$ and factorise its numerator.
\item Draw up the variation table of $C$ on $]0,+\infty[$.
\item Find the value of $x$ that minimises the average cost and calculate this minimum.
\item Sketch a consistent graph of the curve.
\end{enumerate}
\end{exercise}
\begin{solution}\textbf{Marking scheme:} $6+6+8$. \textbf{Solution:} $D_f=\R\setminus\{0\}$, $f'(x)=-\dfrac{16}{x^2}<0$, asymptotes $x=0$ and $y=2$. For the quotient, order $-2,0,4$ and exclude $0$. For the cost, $C'(x)=1-\dfrac{16}{x^2}=\dfrac{(x-4)(x+4)}{x^2}$; minimum at $x=4$, with value $20$.\end{solution}

\mxtest{4}{B}{Reciprocal function}{calculate, reason, represent, communicate}
\begin{exercise}[1 -- Reciprocal function]\mxpoints{6}
Consider $f(x)=2+\dfrac{9}x$.
\begin{enumerate}
\item Give its domain.
\item Calculate its derivative and study its sign.
\item Give its asymptotes and its direction of variation on each interval of the domain.
\end{enumerate}
\end{exercise}
\begin{exercise}[2 -- Sign table]\mxpoints{6}
Solve $\dfrac{(x-3)(x+2)}x\geq0$. The table must show the zeros and the excluded value.
\end{exercise}
\begin{exercise}[3 -- Average cost]\mxpoints{8}
For $x>0$, an average cost is modelled by $C(x)=x+\dfrac{9}x+12$.
\begin{enumerate}
\item Calculate $C'(x)$ and factorise its numerator.
\item Draw up the variation table of $C$ on $]0,+\infty[$.
\item Find the value of $x$ that minimises the average cost and calculate this minimum.
\item Sketch a consistent graph of the curve.
\end{enumerate}
\end{exercise}
\begin{solution}\textbf{Marking scheme:} $6+6+8$. \textbf{Solution:} $D_f=\R\setminus\{0\}$, $f'(x)=-\dfrac{9}{x^2}<0$, asymptotes $x=0$ and $y=2$. For the quotient, order $-2,0,3$ and exclude $0$. For the cost, $C'(x)=1-\dfrac{9}{x^2}=\dfrac{(x-3)(x+3)}{x^2}$; minimum at $x=3$, with value $18$.\end{solution}

\section*{Chapter 5 -- Statistics with two quantitative variables}
\mxtest{5}{A}{Two-variable statistics}{represent, model, calculate, interpret}
\begin{exercise}[1 -- Constructing and describing]\mxpoints{6}
We observe the pairs $(1;8)$, $(2;11)$, $(3;13)$, $(4;17)$, $(5;19)$ and $(6;23)$.
\begin{enumerate}
\item Draw the scatter plot on a correctly labelled coordinate system.
\item Calculate the coordinates of the mean point.
\item Describe the trend and the spread without concluding that there is causality.
\end{enumerate}
\end{exercise}
\begin{exercise}[2 -- Using a linear regression]\mxpoints{6}
A linear regression line is $y=3.2x+12$, established for $x\in[2,15]$.
\begin{enumerate}
\item Estimate $y$ for $x=10$ and characterise the estimate.
\item Estimate $x$ when $y=60$.
\item Estimate $y$ for $x=40$, then explain why this result must be treated with caution.
\end{enumerate}
\end{exercise}
\begin{exercise}[3 -- Change of variable]\mxpoints{8}
Let $z=\log(y)$, and the linear regression of $(x,z)$ is $z=0.06x+1.1$.
\begin{enumerate}
\item Express $y$ in terms of $x$.
\item Estimate $y$ for $x=8$.
\item Explain the benefit of the change of variable and one limitation of extrapolation.
\end{enumerate}
\end{exercise}
\begin{solution}\textbf{Marking scheme:} $6+6+8$. \textbf{Solution:} Calculate $\bar x$ and $\bar y$ from the six values. With the line, substitute $x$ or solve an equation; $x=10$ is an interpolation and $x=40$ a distant extrapolation. The final model is $y=10^{0.06x+1.1}$.\end{solution}

\mxtest{5}{B}{Two-variable statistics}{represent, model, calculate, interpret}
\begin{exercise}[1 -- Constructing and describing]\mxpoints{6}
We observe the pairs $(1;8)$, $(2;11)$, $(3;13)$, $(4;17)$, $(5;22)$ and $(6;23)$.
\begin{enumerate}
\item Draw the scatter plot on a correctly labelled coordinate system.
\item Calculate the coordinates of the mean point.
\item Describe the trend and the spread without concluding that there is causality.
\end{enumerate}
\end{exercise}
\begin{exercise}[2 -- Using a linear regression]\mxpoints{6}
A linear regression line is $y=2.4x+18$, established for $x\in[2,15]$.
\begin{enumerate}
\item Estimate $y$ for $x=10$ and characterise the estimate.
\item Estimate $x$ when $y=60$.
\item Estimate $y$ for $x=40$, then explain why this result must be treated with caution.
\end{enumerate}
\end{exercise}
\begin{exercise}[3 -- Change of variable]\mxpoints{8}
Let $z=\log(y)$, and the linear regression of $(x,z)$ is $z=0.06x+1.1$.
\begin{enumerate}
\item Express $y$ in terms of $x$.
\item Estimate $y$ for $x=8$.
\item Explain the benefit of the change of variable and one limitation of extrapolation.
\end{enumerate}
\end{exercise}
\begin{solution}\textbf{Marking scheme:} $6+6+8$. \textbf{Solution:} Calculate $\bar x$ and $\bar y$ from the six values. With the line, substitute $x$ or solve an equation; $x=10$ is an interpolation and $x=40$ a distant extrapolation. The final model is $y=10^{0.06x+1.1}$.\end{solution}

\section*{Chapter 6 -- Conditional probability}
\mxtest{6}{A}{Conditional probability}{model, calculate, reason, interpret}
\begin{exercise}[1 -- Weighted tree]\mxpoints{6}
A proportion 0.35 of orders are express. The late-delivery rate is 0.05 among express orders and 0.12 among the others.
\begin{enumerate}
\item Define the events $E$ and $R$, then build the complete weighted tree.
\item Calculate $P(E\cap R)$.
\item Calculate $P(R)$ using the law of total probability.
\item Calculate $P_R(E)$ and interpret.
\end{enumerate}
\end{exercise}
\begin{exercise}[2 -- Independence]\mxpoints{6}
In another study, $P(A)=0.4$, $P(B)=0.25$ and $P(A\cap B)=0.10$.
\begin{enumerate}
\item Calculate $P(A)P(B)$.
\item Are the events independent? Justify.
\item Calculate $P_A(B)$ and confirm the conclusion.
\end{enumerate}
\end{exercise}
\begin{exercise}[3 -- Frequency table]\mxpoints{8}
Out of $800$ customers, $320$ have a premium card; $96$ premium customers and $72$ non-premium customers cancelled. Build the complete table, then calculate the conditional probability of being premium given that the customer cancelled.
\end{exercise}
\begin{solution}\textbf{Marking scheme:} $8+6+6$. \textbf{Solution:} $P(R)=P(E)P_E(R)+P(\overline E)P_{\overline E}(R)$ then $P_R(E)=\dfrac{P(E\cap R)}{P(R)}$. For independence, compare the intersection with the product $0.10$. In the table, there are $168$ cancellations, so $P_R(P)=\dfrac{96}{168}=\dfrac47$.\end{solution}

\mxtest{6}{B}{Conditional probability}{model, calculate, reason, interpret}
\begin{exercise}[1 -- Weighted tree]\mxpoints{6}
A proportion 0.45 of orders are express. The late-delivery rate is 0.08 among express orders and 0.15 among the others.
\begin{enumerate}
\item Define the events $E$ and $R$, then build the complete weighted tree.
\item Calculate $P(E\cap R)$.
\item Calculate $P(R)$ using the law of total probability.
\item Calculate $P_R(E)$ and interpret.
\end{enumerate}
\end{exercise}
\begin{exercise}[2 -- Independence]\mxpoints{6}
In another study, $P(A)=0.4$, $P(B)=0.25$ and $P(A\cap B)=0.14$.
\begin{enumerate}
\item Calculate $P(A)P(B)$.
\item Are the events independent? Justify.
\item Calculate $P_A(B)$ and confirm the conclusion.
\end{enumerate}
\end{exercise}
\begin{exercise}[3 -- Frequency table]\mxpoints{8}
Out of $800$ customers, $320$ have a premium card; $96$ premium customers and $72$ non-premium customers cancelled. Build the complete table, then calculate the conditional probability of being premium given that the customer cancelled.
\end{exercise}
\begin{solution}\textbf{Marking scheme:} $8+6+6$. \textbf{Solution:} $P(R)=P(E)P_E(R)+P(\overline E)P_{\overline E}(R)$ then $P_R(E)=\dfrac{P(E\cap R)}{P(R)}$. For independence, compare the intersection with the product $0.10$. In the table, there are $168$ cancellations, so $P_R(P)=\dfrac{96}{168}=\dfrac47$.\end{solution}

\section*{Chapter 7 -- Discrete random variables and the binomial distribution}
\mxtest{7}{A}{Random variables and the binomial distribution}{model, calculate, reason, interpret}
\begin{exercise}[1 -- Expected value of a game]\mxpoints{6}
A game costs $3$ euros. The gross winnings are $0$, $5$ or $15$ euros with probabilities $0.65$, $0.25$ and $0.10$.
\begin{enumerate}
\item Build the distribution of the net gain $G$.
\item Calculate $E(G)$.
\item Is the game favourable to the player? Interpret without claiming what will happen in a single play.
\end{enumerate}
\end{exercise}
\begin{exercise}[2 -- Binomial distribution]\mxpoints{6}
A customer responds favourably with probability $0.30$. An adviser contacts $10$ customers independently. Let $X$ be the number of favourable responses.
\begin{enumerate}
\item Justify the distribution of $X$ and give its parameters.
\item Calculate $P(X=2)$.
\item Calculate $P(X=0)$ then $P(X\geq1)$.
\item Calculate and interpret $E(X)$.
\end{enumerate}
\end{exercise}
\begin{exercise}[3 -- Pascal's triangle]\mxpoints{8}
\begin{enumerate}
\item Write row $n=6$ of Pascal's triangle.
\item Read off $\binom62$ and $\binom64$.
\item Explain the role of $\binom nk$ in the binomial formula.
\end{enumerate}
\end{exercise}
\begin{solution}\textbf{Marking scheme:} $6+8+6$. \textbf{Solution:} The net gain is $-3,2,12$; $E(G)=-1.95+0.50+1.20=-0.25$ euro. $X\sim\mathcal B(10,0.30)$; apply the formula, use $1-P(X=0)$ for ``at least one'', and $E(X)=np$. Row $n=6$ is $1,6,15,20,15,6,1$.\end{solution}

\mxtest{7}{B}{Random variables and the binomial distribution}{model, calculate, reason, interpret}
\begin{exercise}[1 -- Expected value of a game]\mxpoints{6}
A game costs $3$ euros. The gross winnings are $0$, $5$ or $15$ euros with probabilities $0.65$, $0.25$ and $0.10$.
\begin{enumerate}
\item Build the distribution of the net gain $G$.
\item Calculate $E(G)$.
\item Is the game favourable to the player? Interpret without claiming what will happen in a single play.
\end{enumerate}
\end{exercise}
\begin{exercise}[2 -- Binomial distribution]\mxpoints{6}
A customer responds favourably with probability $0.25$. An adviser contacts $8$ customers independently. Let $X$ be the number of favourable responses.
\begin{enumerate}
\item Justify the distribution of $X$ and give its parameters.
\item Calculate $P(X=2)$.
\item Calculate $P(X=0)$ then $P(X\geq1)$.
\item Calculate and interpret $E(X)$.
\end{enumerate}
\end{exercise}
\begin{exercise}[3 -- Pascal's triangle]\mxpoints{8}
\begin{enumerate}
\item Write row $n=6$ of Pascal's triangle.
\item Read off $\binom62$ and $\binom64$.
\item Explain the role of $\binom nk$ in the binomial formula.
\end{enumerate}
\end{exercise}
\begin{solution}\textbf{Marking scheme:} $6+8+6$. \textbf{Solution:} The net gain is $-3,2,12$; $E(G)=-1.95+0.50+1.20=-0.25$ euro. $X\sim\mathcal B(8,0.25)$; apply the formula, use $1-P(X=0)$ for ``at least one'', and $E(X)=np$. Row $n=6$ is $1,6,15,20,15,6,1$.\end{solution}
